\chapter{Avionics and Airframe}\label{cha:avionics-airframe}

This chapter discusses XCSoar as a subsystem of the aircraft.  It
covers the integration of XCSoar with external devices, including GPS,
switches and sensors, and aircraft radio transceivers and other
devices.  Integration with FLARM is covered in
Chapter~\ref{cha:airspace}, and integration with variometers is
covered in Chapter~\ref{cha:atmosph}.

\section{Battery life}

Most modern PDAs are designed for short sporadic use and so do not
have a very good battery capacity when considering the duration of
cross-country soaring flights.  It is recommended to power the PDA
externally, via a transformer connected to the glider battery.  This
installation should be performed by appropriately qualified personnel,
and should contain a fuse and a manual isolation switch.

The greatest cause of power drain by the PDA is the LCD back-light,
however domestic PDAs are not particularly bright so they may need to
have the back-light up full. However, for EFIS systems such as Altair,
it is recommended to use the lowest back-light settings that are
comfortable.

When operating PDAs under internal battery, XCSoar detects a low
battery condition and allows the operating system to shut down and
preserve the memory.  In addition, it can be set up to blank the
screen after a period of inactivity, so that it can reduce the power
consumption.  When the screen is blanked, pressing one of the hardware
buttons on the PDA activates the screen again.  When a status message
is issued by the system, the screen becomes activated.

Another way to help conserve battery power is to reduce the
computational load by turning off certain features.  Drawing terrain
and long snail trails contribute significantly to the CPU load.

For Altair/Vega systems, the external supply voltage is displayed on
the system status dialog (see Section~\ref{sec:system-status-dialog}).

\section{GPS connection}

XCSoar requires a 3D GPS fix for its navigation functions.

\subsection*{GPS status}

GPS status icons and text may appear on the bottom edge of the
map display to indicate:

\begin{tabular}{c c}%{c c}
\includegraphics[angle=0,width=0.75cm,keepaspectratio='true']{icons/gps_acquiring.pdf} & \includegraphics[angle=0,width=0.75cm,keepaspectratio='true']{icons/gps_disconnected.pdf}\\
(a) & (b)
\end{tabular}

\begin{description}
\item[Waiting for GPS fix (a)]  Better reception
  or additional time to search for satellites is required. The GPS may have a 2D fix.
  The aircraft symbol disappears while there is no 3D fix.
\item[GPS not connected (b)]  No communication with the GPS is received.
  This indicates an error in the comm port settings or the GPS device may
  be disconnected or switched off.
\end{description}

When the GPS is not connected for more than one minute, XCSoar
automatically attempts to restart communication with the device and
will then resume waiting.  This method has shown to provide the most
reliable way of recovering from communication errors.

XCSoar can handle up to two GPS sources and it uses them to provide
redundancy.  This means that if the primary GPS source drops out,
XCSoar will use the GPS data from the second source.  If both sources
have valid fixes, the second source is ignored.  For this reason, it
is recommended to have the GPS source with the best antenna or
reliability as the primary device.

\subsection*{GPS altitude}

Some older GPS units (and some new ones) do not output altitude
relative to mean sea level, rather they output elevation with respect
to the WGS84 ellipsoid.  XCSoar detects when this occurs and applies
the ellipsoid to geoid offset according to an internal tabulated data
at two degree spacing.  This is not required for FLARM units or Altair
Pro, which correctly output MSL altitude.

\section{Switch inputs}

XCSoar supports monitoring of switches and sensors connected to the
host computer, for the purpose of providing situational awareness
feedback, alerts, or as general-purpose user-interface input devices.
Several mechanisms are available for interfacing to switches and
sensors:
\begin{description}
\item[Serial device]  Certain intelligent variometers such as
 triadis engineering's Vega have multiple airframe switches
 and pass this information on to the PDA or EFIS as special
 NMEA sentences.
\item[1-Wire device]  triadis engineering's Altair glide computer
 and Vega variometer provide a 1-Wire peripheral bus to which
 various digital and analog sensors can be attached.
\item[Bluetooth device]  Many Pocket PC devices support wireless
 connection to a Bluetooth Game-Pad device that has several buttons.
 This is more suited to user-interface input devices than airframe
 monitoring.
\end{description}

A custom `input events' file determines how switch and sensor
inputs are processed.

A standard set of airframe inputs are defined as:
\begin{itemize}
\item Airbrake
\item Flap position (positive/landing flap, neutral, negative/reflex)
\item Landing gear
\end{itemize}

This set is expected to expand to include engine and fuel monitoring.

Other logical inputs from Vega include computed quantities relating to
specific airframe alerts and aircraft operating envelope warnings, for
example ``airbrake extended and gear retracted''.  

Refer to the Vega documentation and {\em XCSoar Advanced Configuration
Manual} for more details on switch inputs and how they may be used.

\section{Switch dialog}

A dialog displaying switch states for the Vega variometer
is available from the menu:
\begin{quote}
\bmenu{Config}\blink\bmenu{Config}\blink\bmenu{Vario}\blink\bmenu{Airframe
Switches}
\end{quote}

This dialog is updated in real-time, allowing the pilot
to check the correct functioning of switches during daily
inspection tests or before takeoff. 

\begin{center}
\includegraphics[angle=0,width=0.7\linewidth,keepaspectratio='true']{figures/dialog-switches.png}
\end{center}

\section{Aircraft radio transceiver}

Monitoring and setting the active and standby frequencies of aircraft
transceivers with serial connections are currently in development.
The goal is to provide the functionality whereby the radio frequency
can be set with one button press (requiring confirmation by the pilot)
when the aircraft flies into controlled airspace or enters an
airfield's advisory/mandatory broadcast traffic zone.

\section{Supported variometers}\label{sec:supported-varios}

Currently XCSoar supports inputs from a very long list of intelligent
variometers. To catch up with the growing number of devices and their supported 
features please consult the XCSoar homepage at 
\url{http://www.xcsoar.org/discover/features.html#supported_devices}.

Note that the level of support for each device varies, and not all manufacturers
release their protocols to allow the XCSoar developers to provide full support. 

\section{Slave mode}

A device type in the configuration settings, ``NMEA Out'' is defined
for use in joining two Altair or PDA systems in a master-slave mode.
In the master, the second com device can be set to NMEA Out, and all
data received in the first com device (as well as outgoing data) will
be sent to the slave.  

As an example where two Altairs are being connected together, in the
slave, the first com device can then be set to ``Vega'' or ``Altair
Pro'' and this system receives all data as if it came from the
Master's GPS and connected instruments (Vega, FLARM etc).

\section{Interface to external loggers}

XCSoar has support for declaration and handling special sentences used
by commercial flight loggers.  Devices that are supported but do not
include declaration support include Posigraph and Colibri.  These will
be improved in the future.

\subsection*{Volkslogger}

Uploading of the current task to the Volkslogger is supported.  The
device type in the system settings must be set to `Volkslogger'.

When the \button{Declare} button is pressed in the task editor, the
current task, and pilot name, aircraft type and competition ID will be
uploaded to the Volkslogger.

Note that uploading the task erases the Volkslogger's database of
waypoints!

The support for synchronising the waypoint database between XCSoar and
the Volkslogger will be improved in the future. 

\subsection*{EW logger}

Uploading of the current task to EW loggers is supported.  The
device type in the system settings must be set to `EW'.

When the \button{Declare} button is pressed in the task editor, the
current task, will be uploaded to the Volkslogger.

\subsection*{CAI 302}

Uploading of the current task to the CAI 302 is supported.  The device
type in the system settings must be set to `CAI 302.

When the \button{Declare} button is pressed in the task editor, the
current task, and pilot name, aircraft type and competition ID will be
uploaded to the CAI 302.


\section{System status dialog}\label{sec:system-status-dialog}

The system status dialog (see Section~\ref{sec:dialog-windows}) is
used primarily as a systems check, to see how the host computer and
connected devices are performing.

This is accessed via the menu under 
\begin{quote}
\bmenu{Info}\blink\bmenu{Info}\blink\bmenu{Status}
\end{quote}
and then selecting the page `System'.

All dynamic values (e.g.\ battery voltage, number of satellites in
view) are updated continuously.
